(1) 每层电子速度扰动$v_j$满足$-i\omega mv_j=-eE_{xj}=+ieq\phi_j$，因此$v_j=-eq\phi_j/(m\omega)$。电子连续性$-i\omega\delta n_j+iqn_sv_j=0$给$\delta n_j=qn_sv_j/\omega=-n_se q^2\phi_j/(m\omega^2)$。面电荷$\delta\sigma_j=-e\delta n_j=n_se^2q^2\phi_j/(m\omega^2)$，同电势扰动正相关，符合电子负电荷双重符号。
(2) 二维傅里叶空间内，Poisson方程$(\partial_z^2-q^2)\phi=-\sum_j\sigma_j\delta(z-z_j)/\epsilon$的Green函数是$e^{-q|z-z_j|}/(2\epsilon q)$。在两层位置求值，$\binom{\phi_1}{\phi_2}=(2\epsilon q)^{-1}\begin{pmatrix}1&e^{-qd}\\e^{-qd}&1\end{pmatrix}\binom{\sigma_1}{\sigma_2}$。它的正交本征向量$(1,1)$与$(1,-1)$分别描述同相和反相电荷涨落。
(3) 联立电荷响应$\sigma_j=[n_se^2q^2/(m\omega^2)]\phi_j$，令矩阵行列式为零，得到$\omega_\pm^2(q)=\omega_{2D}^2(q)[1\pm e^{-qd}]$，其中单层二维等离激元$\omega_{2D}^2(q)=n_se^2q/(2m\epsilon)$。因$1\pm e^{-qd}>0$，理想无耗散谱均为实数，且对称支频率高于反对称支。
(4) 当$qd\ll1$，$e^{-qd}=1-qd+O(q^2d^2)$，故$\omega_+^2\simeq n_se^2q/(m\epsilon)$仍为$\sqrt q$色散，而$\omega_-^2\simeq n_se^2dq^2/(2m\epsilon)$，所以反相模式变成声学型$\omega_-\simeq q\sqrt{n_se^2d/(2m\epsilon)}$。当$qd\gg1$，层间库仑耦合指数消失，两支均趋单层$\omega_{2D}(q)$。有限层厚与电子压强可在更短波长修改结果。
